\section{Ánh xạ Kiểu dữ liệu và Ràng buộc (Đầu vào)}

Theo quy trình xử lý, đầu vào của thiết kế vật lý bao gồm Lược đồ logic đã được chuẩn hóa (đạt chuẩn BCNF ở Chương 3), cùng với tập hợp các khóa và ràng buộc. Hệ thống OctaPoint tiến hành ánh xạ các yếu tố này sang hệ quản trị SQL Server như sau:

\subsection{Ánh xạ Kiểu dữ liệu (Data Type Mapping)}
\begin{itemize}
    \item \textbf{Khóa định danh (ID):} Sử dụng \texttt{VARCHAR(36)} để lưu UUID thay vì số nguyên tự tăng (\texttt{INT IDENTITY}) nhằm hỗ trợ sinh khóa phân tán và bảo mật.
    \item \textbf{Dữ liệu tài chính:} Các cột \texttt{Amount}, \texttt{AvailableBalance} bắt buộc dùng kiểu \texttt{DECIMAL(18,2)} để tránh sai số dấu phẩy động trong quá trình cộng/trừ điểm thưởng.
    \item \textbf{Dữ liệu phi cấu trúc:} Cột \texttt{SystemConfig} (bảng TENANT) sử dụng kiểu \texttt{NVARCHAR(MAX)} để lưu trữ cấu trúc JSON, tận dụng khả năng phân tích chuỗi động tích hợp sẵn của SQL Server.
\end{itemize}

\subsection{Ràng buộc Toàn vẹn (Constraints)}
\begin{itemize}
    \item \textbf{Khóa chính (Primary Key - PK):} Bắt buộc tồn tại trên mọi bảng để xác định duy nhất một bản ghi trong khối dữ liệu (SQL Server sẽ tự động tạo Clustered Index trên cột này).
    \item \textbf{Khóa ngoại (Foreign Key - FK):} Thiết lập ràng buộc tham chiếu chặt chẽ (mặc định \texttt{ON DELETE NO ACTION}) để ngăn chặn việc tạo ra dữ liệu mồ côi (orphan data) khi xóa các thực thể cha như Merchant hay Campaign.
    \item \textbf{Ràng buộc Unique (Tính duy nhất):} Áp dụng cho \texttt{Email} (bảng USER), \texttt{SuiTxDigest} (bảng CREDIT\_TRANSACTION) để chống trùng lặp. Đặc biệt, thiết lập ràng buộc \texttt{UNIQUE(CustomerID, TenantID)} trên bảng \texttt{CREDIT} để đảm bảo tuân thủ tuyệt đối quy tắc nghiệp vụ: Mỗi khách hàng chỉ có duy nhất 1 ví điểm tại 1 chuỗi bán lẻ.
\end{itemize}